\begin{lemma}[$n$-dimensional Hermite reproduction identity]
Fix $n \in \mathbb{N}$ and let $a \in \mathbb{R}$ satisfy $0 \leq a \leq 1$.  Write
$\tau = \sqrt{1-a^2}$.  Then for every multi-index $\alpha \in \mathbb{N}_0^n$ and
every $x \in \mathbb{R}^n$,
\[
  \int_{\mathbb{R}^n} h_\alpha\bigl(a\,x + \tau\,y\bigr) \, d\gamma_n(y)
    = a^{|\alpha|}\, h_\alpha(x),
\]
where $h_\alpha$ is the $n$-dimensional Hermite function of \noderef{HermiteFun},
$|\alpha| = \alpha_1 + \cdots + \alpha_n$ is the order of \noderef{MultiIndexSum},
and $\gamma_n$ is the Gaussian measure of \noderef{GaussMeasure}.

The integral on the left is taken against the \emph{Gaussian} measure $\gamma_n$,
which already carries the normalising factor $\pi^{-n/2}$ and the weight
$e^{-|y|^2}$ in its density; no further weight appears in the integrand.  This is
the direct multi-index companion of \noderef{Hermite1Reproduction}, whose
one-dimensional statement is instead written as a Lebesgue integral $d\sigma$ with
the weight $\pi^{-1/2}e^{-\sigma^2}$ displayed in the integrand.  The two
conventions describe the same object in dimension one, and the proof below passes
between them explicitly.
\end{lemma}
\begin{proof}
\emph{Step 1: unfold $\gamma_n$ as a weighted Lebesgue integral.}
By \noderef{GaussMeasure}, $\gamma_n$ is Lebesgue measure on $\mathbb{R}^n$ with
density $\rho(y) = \pi^{-n/2}e^{-|y|^2}$.  That density is measurable and finite
at every point, so integration against $\gamma_n$ is integration of the
density-weighted integrand against Lebesgue measure:
\[
  \int_{\mathbb{R}^n} h_\alpha(a x + \tau y)\, d\gamma_n(y)
    = \int_{\mathbb{R}^n} \pi^{-n/2} e^{-|y|^2}\, h_\alpha(a x + \tau y)\, dy .
\]

\emph{Step 2: the integrand factorizes over the coordinates.}
Fix $y \in \mathbb{R}^n$.  Three factorizations combine.  First,
$|y|^2 = \sum_{j=1}^n y_j^2$, so
$e^{-|y|^2} = \prod_{j=1}^n e^{-y_j^2}$.  Second, the constant splits as
$\pi^{-n/2} = \prod_{j=1}^n \pi^{-1/2}$, since the product of $n$ equal factors
$\pi^{-1/2}$ is $\pi^{-n/2}$.  Third, by \noderef{HermiteFun} the
$n$-dimensional Hermite function is the product of one-dimensional ones over the
coordinates, and the $j$th coordinate of $a x + \tau y$ is $a x_j + \tau y_j$,
so
\[
  h_\alpha(a x + \tau y)
    = \prod_{j=1}^n h_{\alpha_j}\bigl(a x_j + \tau y_j\bigr)
\]
with $h_k$ the one-dimensional Hermite function of \noderef{Hermite1}.
Multiplying the three and regrouping the factors coordinatewise,
\begin{equation*}
  \pi^{-n/2} e^{-|y|^2}\, h_\alpha(a x + \tau y)
    = \prod_{j=1}^n \Bigl( \pi^{-1/2} e^{-y_j^2}\,
        h_{\alpha_j}\bigl(a x_j + \tau y_j\bigr) \Bigr) . \tag{1}
\end{equation*}

\emph{Step 3: the integral of a product splits.}
Lebesgue measure on $\mathbb{R}^n$ is the $n$-fold product of Lebesgue measure on
$\mathbb{R}$, the identification of $\mathbb{R}^n$ with the product being measure
preserving.  Each factor in $(1)$ depends on one coordinate only, so Fubini's
theorem for a finite product of measures applies and gives
\[
  \int_{\mathbb{R}^n} \prod_{j=1}^n \Bigl( \pi^{-1/2} e^{-y_j^2}\,
      h_{\alpha_j}(a x_j + \tau y_j)\Bigr) dy
  = \prod_{j=1}^n \int_{\mathbb{R}} \pi^{-1/2} e^{-\sigma^2}\,
      h_{\alpha_j}(a x_j + \tau \sigma)\, d\sigma .
\]

\emph{Step 4: evaluate each one-dimensional factor.}
The hypotheses $0 \leq a \leq 1$ are exactly what
\noderef{Hermite1Reproduction} requires of its contraction parameter, and
$\tau = \sqrt{1-a^2}$ is the spreading factor that lemma pairs with it.  Applying
it with $r = a$, $k = \alpha_j$ and $t = x_j$ evaluates the $j$th factor as
\[
  \int_{\mathbb{R}} \pi^{-1/2} e^{-\sigma^2}\,
      h_{\alpha_j}(a x_j + \tau \sigma)\, d\sigma
    = a^{\alpha_j}\, h_{\alpha_j}(x_j) .
\]

\emph{Step 5: multiply.}
Taking the product over $j$ and separating the two kinds of factor,
\[
  \prod_{j=1}^n a^{\alpha_j}\, h_{\alpha_j}(x_j)
    = \Bigl(\prod_{j=1}^n a^{\alpha_j}\Bigr)
      \prod_{j=1}^n h_{\alpha_j}(x_j)
    = a^{\sum_{j} \alpha_j}\, h_\alpha(x)
    = a^{|\alpha|}\, h_\alpha(x),
\]
where the exponents were collected using $\prod_j a^{\alpha_j} = a^{\sum_j
\alpha_j}$, the remaining product was recognised as $h_\alpha(x)$ by
\noderef{HermiteFun}, and $\sum_j \alpha_j = |\alpha|$ by
\noderef{MultiIndexSum}.  Chaining Steps 1 through 5 gives the assertion.
\end{proof}
